\chapter{Introduction}\label{ch:introduction}

A single rendered frame in an interactive vector map has to fit a tile fetch, an evaluation of style expressions over every visible feature, and a \ac{GPU} submit inside the same frame budget, and on a slow connection or a battery-bound device any of the three can become the bottleneck.
MapLibre GL~JS\footnote{We target MapLibre GL~JS, the open-source fork of Mapbox GL~JS, because its source code permits the static analysis and runtime instrumentation that a style optimizer requires.
The proprietary Mapbox codebase does not allow independent research evaluation.} and Mapbox GL~JS deliver these maps by streaming tiled geometry to the browser and letting a client-side style document drive a \ac{WebGL}- or \ac{WebGPU}-based rendering pipeline.
This architecture is what makes panning, zooming, and rotation smooth when it works, and what makes them stutter when it does not.

As map deployments scale - from personal projects to platforms serving millions of daily users - the cost of suboptimal performance compounds.
Larger tiles increase transfer time and bandwidth cost, redundant style layers waste rendering cycles and unnecessary properties inflate decoding overhead.
On mobile devices, these inefficiencies translate directly into battery drain and thermal throttling, though we leave direct energy measurement outside the scope of this work.
On unstable network connections, such as on trains or aeroplanes, these translate to degraded user experiences.

Despite this, the map rendering and geospatial community has largely treated style authoring and data preparation as independent concerns: styles are hand-tuned for visual appearance, while tiles are encoded with general-purpose content and compression.
The interaction between the two is left unexploited.
For example, a style filter that matches on a property value may, after constant folding with tile-level statistics, simplify to a form that no longer references this property.
This in turn allows the data encoder to drop the corresponding column from every tile - a saving that neither style simplification nor data pruning could achieve independently.
We ask whether such cross-level interactions produce measurably super-additive gains, or whether the levels compose merely additively, as a central empirical question of this work that we formalise in \cref{sec:problem_statement}.

\section{Research Goals}\label{sec:research_goals}

We investigate automatic, style- and data-driven co-optimization for MapLibre-based map rendering pipelines.
Three research questions guide our work:

\begin{enumerate}[label=\textbf{R\arabic*}]
    \item To what extent \textbf{can automatic data and style co-optimization improve rendering performance and reduce map tile size}?
    \item \textbf{Which techniques} - such as expression simplification, dead layer elimination, layer merging, tile shaving, and columnar re-encoding - \textbf{are most effective}, \textbf{and how do they interact}?  In particular, R2 asks whether the composition of style-level and data-level optimizations is additive or synergistic.
    \item \textbf{How expensive are these optimizations} in terms of preprocessing time, and what are the \textbf{tradeoffs between decoding latency, rendering throughput, and transfer size}?
\end{enumerate}

To answer these questions, we design and implement a standalone optimizer that consumes a MapLibre style document and optional tile-level statistics, and produces a equivalent (\cref{eq:visual_equivalence}) style whose rendering performance and network utilsation is measurably improved.

\newpage
\section{Contributions}\label{sec:contributions}

We make the following main contributions:

\begin{enumerate}
    \item A \textbf{three-level optimization pipeline} comprising expression-level rewrites (constant folding, algebraic simplification, stats-driven folding, selectivity reordering), structural passes (dead layer elimination, metadata refinement, layer merging), and data-level transformations (tile shaving, encoding strategy selection), modeled on peephole optimizers and query planners from the compiler and database literature.
    We decompose the pipeline along three distinct input representations:
    Expression passes operate on raw \ac{JSON} arrays with no schema knowledge, enabling fast, composable peephole rules.
    Structural passes require a typed model of the full style document to reason across layers and sources.
    And data passes operate on tile content rather than the style, requiring a different input representation entirely.

    We treat the three levels as complementary rather than co-equal:
    each targets a different performance bottleneck, and their relative importance varies by deployment scenario.

    \item \textbf{Style-driven tile shaving}, the dominant source of load-time improvements we observe in this work, which derives a pruning advisory from the optimized style to eliminate unused properties, geometry types, property values, and feature IDs from tile data.
    Style-level passes further refine the advisory, though we find the interaction remains additive rather than synergistic (\cref{sub:mlt_discussion}).

	\item An \textbf{automatic encoding strategy selection} for the \ac{MLT} columnar format that reformulates tile encoding as a multi-strategy competition over sort orderings, integer encodings, and string compression strategies
	The approach adapts to the altered data characteristics introduced by style-driven shaving, including changes in value distributions and column cardinalities.

    \item An \textbf{empirical evaluation} across 15 styles, 18 scenarios, and 19~cumulative ablation steps plus one plain-tile-shaving comparison (fixed pass ordering; we discuss interaction effects in \cref{sec:threats}), which we validated end-to-end on ten deep-corpus styles with tile data, demonstrating measurable improvements in tile volume, load time, and rendering stability.
\end{enumerate}

\section{Outline}\label{sec:outline}

\cref{ch:related_work} shows that prior work attacks each level in isolation - tile-data pruners are style-blind, columnar formats ignore which properties the style actually reads, and style-level optimizers are limited to default-stripping analogues - and \cref{ch:theoretical_background} assembles the compiler, query-optimization, and columnar-encoding foundations the pipeline reuses.

\cref{ch:methods} states the joint problem formally in \cref{sec:problem_statement} - producing $(S', T')$ from $(S, T)$ under the pixel-equivalence constraint of \cref{eq:visual_equivalence} - and details the three-level pipeline that solves it: expression rewrites over raw \ac{JSON}, structural passes over a typed style model, and data-level passes that emit a style-driven pruning advisory and re-encode in \ac{MLT} with automatic strategy selection.

\cref{ch:experiments} measures the isolated and combined effects through a cumulative ablation across 15~styles and 18~scenarios; the full pipeline achieves load-time reductions of \qtyrange{63}{74}{\percent} (median \qty{71}{\percent}), \ac{FPS} gains of \qtyrange{100}{200}{\percent}, and a \qty{28}{\percent} tile-volume reduction over the reference \ac{MVT} baseline, but the composition turns out to be \emph{additive rather than synergistic} at the aggregate level.
\cref{ch:conclusion} revisits R1-R3 in \cref{sec:rq_revisited}, names tile shaving and columnar re-encoding as the primary levers, and points to selective-column decoding, learned encoding selection, and online deployment as the natural next steps.
